\documentclass[12pt,a4paper]{article}
\usepackage[utf8]{inputenc}
\usepackage[T1]{fontenc}
\usepackage{amsmath,amsfonts,amssymb}
\usepackage{graphicx}
\usepackage{geometry}
\usepackage{booktabs}
\usepackage{array}
\usepackage{multirow}
\usepackage{float}
\usepackage[hidelinks]{hyperref}
\usepackage{cleveref}
\usepackage{caption}
\usepackage{subcaption}
\usepackage{enumitem}

\geometry{margin=2.5cm}

\crefname{equation}{Eq.}{Eqs.}
\Crefname{equation}{Equation}{Equations}

\title{\textbf{Mathematical Formulation of Energy Market Operations:\\
Renewable Energy Communities and Supplier Competition Analysis}}
\author{Energy Systems Analysis}
\date{\today}

\begin{document}

\maketitle
\tableofcontents
\newpage

%%=============================================================================
\section{Introduction}
%%=============================================================================

This document provides the complete mathematical formulation for analysing
energy market operations under five scenarios that span single- and
multi-supplier structures, the presence or absence of a Renewable Energy
Community (REC), and centrally optimised battery storage.  The simulation
period covers one full calendar year (\textit{2016-01-01} to
\textit{2016-12-31}) at 15-minute resolution, yielding $T = 35{,}136$
settlement intervals.

The physical case study is a nine-node low-voltage distribution community in
rural Carinthia, Austria, with three prosumers (nodes 2, 6, 8) and six
consumers (nodes 1, 3, 4, 5, 7, 9), mapped to the SimBench reference network
\textit{1-LV-rural3--2-no\_sw}.  The total installed PV capacity is
24.48~kWp and the total annual community load is approximately 63\,900~kWh.

%%=============================================================================
\section{Notation}
%%=============================================================================

\subsection{Index Sets}

\begin{description}[leftmargin=4em, labelwidth=3.5em]
  \item[$\mathcal{T}$] Set of settlement intervals, $t \in \{1, \ldots, T\}$,
        $T = 35{,}136$ (quarter-hours in 2016)
  \item[$\mathcal{N}$] Set of all nine community participants (nodes)
  \item[$\mathcal{P}$] Set of prosumers $\mathcal{P} \subset \mathcal{N}$,
        $|\mathcal{P}| = 3$ (nodes 2, 6, 8)
  \item[$\mathcal{C}$] Set of consumers $\mathcal{C} \subset \mathcal{N}$,
        $|\mathcal{C}| = 6$ (nodes 1, 3, 4, 5, 7, 9)
  \item[$\mathcal{S}$] Set of energy suppliers
  \item[$\mathcal{G}_{s}$] Set of balancing groups belonging to supplier
        $s \in \mathcal{S}$
  \item[$\mathcal{R}$] Set of RECs; $|\mathcal{R}| \in \{0, 1\}$ depending on
        scenario
  \item[$\mathcal{M}_{r}$] Set of participants belonging to REC
        $r \in \mathcal{R}$
\end{description}

\subsection{Time Resolution}

Each interval has duration $\Delta t = 0.25$~h.
A quantity $Q_{(\cdot)}^{t}$ expressed in \textit{kW} at interval $t$
corresponds to an energy volume of $Q_{(\cdot)}^{t} \cdot \Delta t$~kWh.
For brevity, all power quantities are given in kW and all energy quantities in
kWh throughout; the factor $\Delta t$ is therefore omitted from the flow
equations where both sides are expressed consistently.

\subsection{Parameters}

\begin{description}[leftmargin=4em, labelwidth=3.5em]
  \item[$\ell_{n}^{t}$]   Actual metered load of participant $n$ at interval
        $t$~(kWh)
  \item[$g_{n}^{t}$]      Actual metered generation of prosumer $n$ at
        interval $t$~(kWh)
  \item[$\hat{\ell}_{n,\mathrm{DA}}^{t}$] Day-ahead load forecast for
        participant $n$
  \item[$\hat{g}_{n,\mathrm{DA}}^{t}$]   Day-ahead generation forecast for
        prosumer $n$
  \item[$\hat{\ell}_{n,\mathrm{ID}}^{t}$] Intra-day load forecast for
        participant $n$
  \item[$\hat{g}_{n,\mathrm{ID}}^{t}$]   Intra-day generation forecast for
        prosumer $n$
  \item[$\pi_{\mathrm{DA}}^{t}$] Day-ahead market clearing price
        (EUR/kWh)
  \item[$\pi_{\mathrm{ID}}^{t}$] Intra-day market price (EUR/kWh)
  \item[$\pi_{\mathrm{imb}}^{t}$] System imbalance price (EUR/kWh)
  \item[$\pi_{\mathrm{ret},s}^{t}$] Retail tariff of supplier $s$ (EUR/kWh)
  \item[$\pi_{\mathrm{fi},s}^{t}$]  Feed-in tariff of supplier $s$ (EUR/kWh)
  \item[$E_{\max}$] Battery usable capacity (kWh); $E_{\max} = 200$~kWh in
        scenario~C3
  \item[$P_{\mathrm{c}}^{\max}$, $P_{\mathrm{d}}^{\max}$] Maximum charge /
        discharge power (kW)
  \item[$\eta_{\mathrm{c}}, \eta_{\mathrm{d}}$] Charge / discharge efficiency;
        $\eta_{\mathrm{c}} = \eta_{\mathrm{d}} = 0.95$
  \item[$E_{\min}, E_{0}$] Minimum SoC and initial SoC (kWh)
\end{description}

\subsection{Decision Variables (Battery, Scenario C3 only)}

\begin{description}[leftmargin=4em, labelwidth=3.5em]
  \item[$P_{\mathrm{c}}^{t}$]  Charge power at interval $t$ (kW),
        $P_{\mathrm{c}}^{t} \geq 0$
  \item[$P_{\mathrm{d}}^{t}$]  Discharge power at interval $t$ (kW),
        $P_{\mathrm{d}}^{t} \geq 0$
  \item[$E^{t}$]               State of charge at end of interval $t$ (kWh)
  \item[$u^{t}$]               Binary mode indicator: $u^{t} = 1$ if charging,
        $u^{t} = 0$ if discharging
  \item[$I^{t}$]               Net grid import (kWh), $I^{t} \geq 0$
  \item[$X^{t}$]               Net grid export (kWh), $X^{t} \geq 0$
\end{description}

%%=============================================================================
\section{Scenario Definitions}
%%=============================================================================

All five scenarios share the same nine-node community and the same SimBench
reference network \textit{1-LV-rural3--2-no\_sw}.  They differ only in market
structure, REC activation, and battery presence.

\begin{table}[H]
\centering
\caption{Scenario overview}
\label{tab:scenarios}
\begin{tabular}{@{}llccc@{}}
\toprule
Scenario & Description & Suppliers & REC & Battery\\
\midrule
A1 & Single supplier, no REC          & 1 ($|\mathcal{S}|=1$) & No  & No \\
A2 & Single supplier, with REC        & 1                     & Yes & No \\
B1 & Multiple suppliers, no REC       & 2 ($|\mathcal{S}|=2$) & No  & No \\
B2 & Multiple suppliers, with REC     & 2                     & Yes & No \\
C3 & Single supplier, REC + battery   & 1                     & Yes & Yes\\
\bottomrule
\end{tabular}
\end{table}

\noindent
In scenarios A1 and A2 all participants are served by a single supplier
$s^{*}$ with one balancing group.  In scenarios B1 and B2 participants are
distributed across two suppliers, each managing one balancing group.  Scenario
C3 extends A2 by adding a centralised 200-kWh MILP-optimised battery whose
schedule is determined before the REC settlement and balancing steps.

The sequential market pipeline executed for each scenario is:

\begin{equation}
\text{Step~i} \;\rightarrow\; \text{Step~ii}
\;\xrightarrow{\text{C3 only}}\; \text{Step~ii-b (battery)}
\;\xrightarrow{\text{A2/B2/C3}}\; \text{Step~iii (REC)}
\;\rightarrow\; \text{Step~iv} \;\rightarrow\; \text{Step~v}
\label{eq:pipeline}
\end{equation}

%%=============================================================================
\section{Step~i: Day-Ahead Market}
%%=============================================================================

\subsection{Balancing-Group Aggregation}

For each balancing group $g \in \mathcal{G}_{s}$ of supplier $s$, the
aggregated day-ahead forecasts are

\begin{equation}
L_{g,\mathrm{DA}}^{t}
  = \sum_{n \in \mathcal{N}_{g}} \hat{\ell}_{n,\mathrm{DA}}^{t}, \qquad
G_{g,\mathrm{DA}}^{t}
  = \sum_{n \in \mathcal{P}_{g}} \hat{g}_{n,\mathrm{DA}}^{t}
\label{eq:bg_da_agg}
\end{equation}

where $\mathcal{N}_{g}$ and $\mathcal{P}_{g}$ denote the participants and the
prosumers assigned to group $g$, respectively.

\subsection{Net Position and Market Commitment}

The net position of balancing group $g$ in the day-ahead market is

\begin{equation}
Q_{g,\mathrm{DA}}^{t} = G_{g,\mathrm{DA}}^{t} - L_{g,\mathrm{DA}}^{t}
\label{eq:da_net}
\end{equation}

The signed position is split into a purchase volume and a sale volume:

\begin{align}
Q_{g,\mathrm{DA}}^{t,+} &= \max\!\bigl(-Q_{g,\mathrm{DA}}^{t},\;0\bigr)
  &&\text{(net purchase commitment)} \label{eq:da_buy}\\
Q_{g,\mathrm{DA}}^{t,-} &= \max\!\bigl(\phantom{-}Q_{g,\mathrm{DA}}^{t},\;0\bigr)
  &&\text{(net sale commitment)} \label{eq:da_sell}
\end{align}

so that only one of $Q_{g,\mathrm{DA}}^{t,+}$ and $Q_{g,\mathrm{DA}}^{t,-}$
is non-zero at any $t$.  The associated financial commitments are

\begin{align}
C_{g,\mathrm{DA}}^{t} &= Q_{g,\mathrm{DA}}^{t,+} \cdot \pi_{\mathrm{DA}}^{t}
  \label{eq:da_cost}\\
R_{g,\mathrm{DA}}^{t} &= Q_{g,\mathrm{DA}}^{t,-} \cdot \pi_{\mathrm{DA}}^{t}
  \label{eq:da_rev}
\end{align}

and the net day-ahead settlement per interval is

\begin{equation}
S_{g,\mathrm{DA}}^{t} = R_{g,\mathrm{DA}}^{t} - C_{g,\mathrm{DA}}^{t}
\label{eq:da_net_settlement}
\end{equation}

%%=============================================================================
\section{Step~ii: Intra-Day Market}
%%=============================================================================

\subsection{Updated Aggregation}

Updated intra-day forecasts yield revised balancing-group totals:

\begin{equation}
L_{g,\mathrm{ID}}^{t}
  = \sum_{n \in \mathcal{N}_{g}} \hat{\ell}_{n,\mathrm{ID}}^{t}, \qquad
G_{g,\mathrm{ID}}^{t}
  = \sum_{n \in \mathcal{P}_{g}} \hat{g}_{n,\mathrm{ID}}^{t}
\label{eq:bg_id_agg}
\end{equation}

\subsection{Adjustment Volumes}

The intra-day adjustment relative to the day-ahead commitment is

\begin{align}
\Delta Q_{g}^{t,+}
  &= Q_{g,\mathrm{ID}}^{t,+} - Q_{g,\mathrm{DA}}^{t,+}
  \label{eq:id_adj_buy}\\
\Delta Q_{g}^{t,-}
  &= Q_{g,\mathrm{ID}}^{t,-} - Q_{g,\mathrm{DA}}^{t,-}
  \label{eq:id_adj_sell}
\end{align}

where $Q_{g,\mathrm{ID}}^{t,\pm}$ are defined analogously to
\cref{eq:da_buy,eq:da_sell} using $G_{g,\mathrm{ID}}^{t} - L_{g,\mathrm{ID}}^{t}$.

\subsection{Intra-Day Settlement}

\begin{align}
C_{g,\mathrm{ID}}^{t} &= \max\!\bigl(\Delta Q_{g}^{t,+},\;0\bigr)
  \cdot \pi_{\mathrm{ID}}^{t}
  \label{eq:id_cost}\\
R_{g,\mathrm{ID}}^{t} &= \max\!\bigl(\Delta Q_{g}^{t,-},\;0\bigr)
  \cdot \pi_{\mathrm{ID}}^{t}
  \label{eq:id_rev}
\end{align}

\subsection{Closing Position}

The closing commitment — used as the reference for the balancing market — is

\begin{align}
\bar{Q}_{g}^{t,+} &= Q_{g,\mathrm{ID}}^{t,+} \label{eq:closing_buy}\\
\bar{Q}_{g}^{t,-} &= Q_{g,\mathrm{ID}}^{t,-} \label{eq:closing_sell}\\
\bar{S}_{g}^{t}   &= \bigl(C_{g,\mathrm{DA}}^{t} + C_{g,\mathrm{ID}}^{t}\bigr)
         - \bigl(R_{g,\mathrm{DA}}^{t} + R_{g,\mathrm{ID}}^{t}\bigr)
  \label{eq:closing_cost}
\end{align}

%%=============================================================================
\section{Step~ii-b: Battery Optimisation (Scenario C3 Only)}
%%=============================================================================

After the intra-day step, a MILP problem is solved to determine an optimal
charge/discharge schedule for the 200-kWh community battery.  The optimisation
uses intra-day forecasts as the best available information at decision time.

\subsection{Objective Function}

\begin{equation}
\min_{P_{\mathrm{c}},\,P_{\mathrm{d}},\,E,\,u,\,I,\,X}
\quad
\sum_{t \in \mathcal{T}}\Bigl[
  I^{t} \cdot \pi_{\mathrm{ret},s}^{t}
  - X^{t} \cdot \pi_{\mathrm{fi},s}^{t}
\Bigr]
\label{eq:battery_obj}
\end{equation}

\subsection{State-of-Charge Dynamics}

\begin{equation}
E^{t} =
\begin{cases}
  E_{0} + \eta_{\mathrm{c}}\,P_{\mathrm{c}}^{t}\,\Delta t
       - \dfrac{P_{\mathrm{d}}^{t}\,\Delta t}{\eta_{\mathrm{d}}}
  & t = 1\\[8pt]
  E^{t-1} + \eta_{\mathrm{c}}\,P_{\mathrm{c}}^{t}\,\Delta t
           - \dfrac{P_{\mathrm{d}}^{t}\,\Delta t}{\eta_{\mathrm{d}}}
  & t \geq 2
\end{cases}
\label{eq:soc_dynamics}
\end{equation}

\subsection{Power and SoC Bounds}

\begin{align}
E_{\min} &\leq E^{t} \leq E_{\max}
  &&\forall\, t \in \mathcal{T} \label{eq:soc_bounds}\\
0 &\leq P_{\mathrm{c}}^{t} \leq P_{\mathrm{c}}^{\max}\,u^{t}
  &&\forall\, t \in \mathcal{T} \label{eq:charge_limit}\\
0 &\leq P_{\mathrm{d}}^{t} \leq P_{\mathrm{d}}^{\max}(1 - u^{t})
  &&\forall\, t \in \mathcal{T} \label{eq:discharge_limit}\\
u^{t} &\in \{0,1\}
  &&\forall\, t \in \mathcal{T} \label{eq:binary}
\end{align}

Constraints \cref{eq:charge_limit,eq:discharge_limit} enforce that the battery
cannot charge and discharge simultaneously.

\subsection{Net Power Balance}

Let $\tilde{L}^{t} = \sum_{n \in \mathcal{N}} \hat{\ell}_{n,\mathrm{ID}}^{t}$
and $\tilde{G}^{t} = \sum_{n \in \mathcal{P}} \hat{g}_{n,\mathrm{ID}}^{t}$
be community-level ID forecasts.  The net load including battery is

\begin{equation}
\tilde{N}^{t} = \tilde{L}^{t} - \tilde{G}^{t}
         + P_{\mathrm{c}}^{t} - P_{\mathrm{d}}^{t}
\label{eq:net_with_battery}
\end{equation}

and the grid boundary condition is

\begin{equation}
I^{t} - X^{t} = \tilde{N}^{t}, \qquad I^{t} \geq 0,\quad X^{t} \geq 0
\label{eq:grid_balance}
\end{equation}

\noindent
The MILP is solved with a GLPK/CBC solver via Pyomo before the REC settlement
and balancing steps, and the resulting net battery output
$(P_{\mathrm{d}}^{t} - P_{\mathrm{c}}^{t})$ is applied as an additional
generation signal in the downstream steps.

%%=============================================================================
\section{Step~iii: REC Settlement (Scenarios A2, B2, C3)}
%%=============================================================================

The REC settlement operates on \emph{actual} meter readings.  In scenario~C3
the battery schedule from Step~ii-b is incorporated beforehand so that the
prosumer hosting the battery has its net generation updated accordingly.

\subsection{Community Totals}

For REC $r \in \mathcal{R}$:

\begin{align}
G_{r}^{t} &= \sum_{n \in \mathcal{P} \cap \mathcal{M}_{r}} g_{n}^{t}
\label{eq:rec_gen_total}\\
L_{r}^{t} &= \sum_{n \in \mathcal{M}_{r}} \ell_{n}^{t}
\label{eq:rec_load_total}
\end{align}

\subsection{Shared Energy}

The internally shared (self-consumed) energy quantum is

\begin{equation}
E_{r}^{\mathrm{shared},t} = \min\!\bigl(G_{r}^{t},\; L_{r}^{t}\bigr)
\label{eq:shared_energy}
\end{equation}

\subsection{Corrected Meter Readings}

The REC redistribution proportionally reduces participant meters.

\paragraph{Generation correction.}
For each prosumer $n \in \mathcal{P} \cap \mathcal{M}_{r}$:

\begin{equation}
\tilde{g}_{n}^{t} = g_{n}^{t}
  - \frac{g_{n}^{t}}{G_{r}^{t}}\,E_{r}^{\mathrm{shared},t}
  \qquad \bigl(G_{r}^{t} > 0\bigr)
\label{eq:gen_correction}
\end{equation}

\paragraph{Load correction.}
For each member $n \in \mathcal{M}_{r}$:

\begin{equation}
\tilde{\ell}_{n}^{t} = \ell_{n}^{t}
  - \frac{\ell_{n}^{t}}{L_{r}^{t}}\,E_{r}^{\mathrm{shared},t}
  \qquad \bigl(L_{r}^{t} > 0\bigr)
\label{eq:load_correction}
\end{equation}

Equations \cref{eq:gen_correction,eq:load_correction} are consistent with the
three limiting cases:
\begin{itemize}
\item \textbf{Perfect balance} ($G_{r}^{t} = L_{r}^{t}$): all generation is
  consumed internally; both corrected quantities equal zero.
\item \textbf{Community deficit} ($G_{r}^{t} < L_{r}^{t}$): all generation is
  consumed; $\tilde{g}_{n}^{t} = 0$; residual import is
  $\tilde{\ell}_{n}^{t} = \ell_{n}^{t}\!\left(1 - G_{r}^{t}/L_{r}^{t}\right)$.
\item \textbf{Community surplus} ($G_{r}^{t} > L_{r}^{t}$): all load is
  satisfied; $\tilde{\ell}_{n}^{t} = 0$; residual export is
  $\tilde{g}_{n}^{t} = g_{n}^{t}\!\left(1 - L_{r}^{t}/G_{r}^{t}\right)$.
\end{itemize}

\subsection{Corrected Balancing-Group Net Load}

For each balancing group $g$ containing REC members, the post-REC net
load used in the balancing and billing steps is

\begin{equation}
\tilde{N}_{g}^{t}
  = \sum_{n \in \mathcal{N}_{g} \cap \mathcal{M}_{r}} \tilde{\ell}_{n}^{t}
  - \sum_{n \in \mathcal{P}_{g} \cap \mathcal{M}_{r}} \tilde{g}_{n}^{t}
\label{eq:bg_corrected_net}
\end{equation}

In scenarios without a REC (A1, B1) the corrected metres equal the originals:
$\tilde{\ell}_{n}^{t} = \ell_{n}^{t}$ and
$\tilde{g}_{n}^{t} = g_{n}^{t}$ for all $n \in \mathcal{N}$.

%%=============================================================================
\section{Step~iv: Balancing Market}
%%=============================================================================

\subsection{Actual Balancing-Group Position}

For each group $g$, the realised net position using corrected metres is

\begin{equation}
A_{g}^{t} =
  \sum_{n \in \mathcal{N}_{g}} \tilde{\ell}_{n}^{t}
  - \sum_{n \in \mathcal{P}_{g}} \tilde{g}_{n}^{t}
\label{eq:bg_actual}
\end{equation}

\subsection{Imbalance}

The deviation from the closing intra-day commitment defines the imbalance:

\begin{equation}
\varepsilon_{g}^{t} = A_{g}^{t} - \bigl(\bar{Q}_{g}^{t,+} - \bar{Q}_{g}^{t,-}\bigr)
\label{eq:imbalance}
\end{equation}

A positive $\varepsilon_{g}^{t}$ means the group consumed more than
contracted; a negative value means it over-generated or under-consumed.

\subsection{Balancing Settlement}

\begin{equation}
\sigma_{g}^{t} = \varepsilon_{g}^{t} \cdot \pi_{\mathrm{imb}}^{t}
\label{eq:balancing_settlement}
\end{equation}

Decomposed into penalty and reward components:

\begin{equation}
\Pi_{g}^{t} = \max(-\sigma_{g}^{t},\;0), \qquad
\Rho_{g}^{t} = \max(\;\sigma_{g}^{t},\;0)
\label{eq:penalty_reward}
\end{equation}

%%=============================================================================
\section{Step~v: Supplier Billing}
%%=============================================================================

\subsection{Individual Member Settlement}

For each prosumer $n \in \mathcal{P}$ the corrected net exchange with the
grid is

\begin{equation}
\tilde{m}_{n}^{t} = \tilde{\ell}_{n}^{t} - \tilde{g}_{n}^{t}
\label{eq:prosumer_net}
\end{equation}

The supplier charges for import and credits for export:

\begin{equation}
\text{Bill}_{n}^{t}
  = \max\!\bigl(\tilde{m}_{n}^{t},\;0\bigr)
     \cdot \pi_{\mathrm{ret},s(n)}^{t}
   + \max\!\bigl(-\tilde{m}_{n}^{t},\;0\bigr)
     \cdot \pi_{\mathrm{fi},s(n)}^{t}
\label{eq:prosumer_bill}
\end{equation}

where $s(n)$ denotes the supplier of participant $n$.  For a consumer
($n \in \mathcal{C}$) there is no generation, so
$\text{Bill}_{n}^{t} = \tilde{\ell}_{n}^{t} \cdot \pi_{\mathrm{ret},s(n)}^{t}$.

\subsection{Retail Revenue and Feed-in Cost of Supplier}

\begin{align}
R_{s}^{\mathrm{ret},t}
  &= \sum_{n \in \mathcal{N}_{s}}
     \max\!\bigl(\tilde{m}_{n}^{t},\;0\bigr)
     \cdot \pi_{\mathrm{ret},s}^{t}
\label{eq:retail_revenue}\\
C_{s}^{\mathrm{fi},t}
  &= \sum_{n \in \mathcal{P}_{s}}
     \max\!\bigl(-\tilde{m}_{n}^{t},\;0\bigr)
     \cdot \pi_{\mathrm{fi},s}^{t}
\label{eq:feedin_cost}
\end{align}

%%=============================================================================
\section{Profit and Loss}
%%=============================================================================

\subsection{Revenue Components}

For supplier $s$ over the full simulation year:

\begin{align}
R_{s}^{\mathrm{mkt}} &= \sum_{t \in \mathcal{T}}\;\sum_{g \in \mathcal{G}_{s}}
  \Bigl(R_{g,\mathrm{DA}}^{t} + R_{g,\mathrm{ID}}^{t}\Bigr)
  \label{eq:rev_mkt}\\
R_{s}^{\mathrm{bal}} &= \sum_{t \in \mathcal{T}}\;\sum_{g \in \mathcal{G}_{s}}
  \Rho_{g}^{t}
  \label{eq:rev_bal}\\
R_{s}^{\mathrm{ret}} &= \sum_{t \in \mathcal{T}} R_{s}^{\mathrm{ret},t}
  \label{eq:rev_ret}\\
R_{s}^{\mathrm{total}} &= R_{s}^{\mathrm{mkt}} + R_{s}^{\mathrm{bal}}
  + R_{s}^{\mathrm{ret}}
  \label{eq:rev_total}
\end{align}

\subsection{Cost Components}

\begin{align}
C_{s}^{\mathrm{mkt}} &= \sum_{t \in \mathcal{T}}\;\sum_{g \in \mathcal{G}_{s}}
  \Bigl(C_{g,\mathrm{DA}}^{t} + C_{g,\mathrm{ID}}^{t}\Bigr)
  \label{eq:cost_mkt}\\
C_{s}^{\mathrm{bal}} &= \sum_{t \in \mathcal{T}}\;\sum_{g \in \mathcal{G}_{s}}
  \Pi_{g}^{t}
  \label{eq:cost_bal}\\
C_{s}^{\mathrm{fi}}  &= \sum_{t \in \mathcal{T}} C_{s}^{\mathrm{fi},t}
  \label{eq:cost_fi}\\
C_{s}^{\mathrm{total}} &= C_{s}^{\mathrm{mkt}} + C_{s}^{\mathrm{bal}}
  + C_{s}^{\mathrm{fi}}
  \label{eq:cost_total}
\end{align}

\subsection{Annual Profit / Loss}

\begin{equation}
\Pi_{s}^{\mathrm{annual}} = R_{s}^{\mathrm{total}} - C_{s}^{\mathrm{total}}
\label{eq:pl}
\end{equation}

%%=============================================================================
\section{Scenario Comparison Framework}
%%=============================================================================

\subsection{REC Impact}

Comparing scenarios with and without the energy community, holding market
structure fixed, isolates the REC effect on supplier cost:

\begin{equation}
\Delta C^{\mathrm{REC}} = C^{\mathrm{total}}_{\text{A2}} - C^{\mathrm{total}}_{\text{A1}}
\label{eq:rec_impact}
\end{equation}

\subsection{Supplier Competition Effect}

Comparing single- against multi-supplier scenarios without REC isolates
market-structure effects:

\begin{equation}
\Delta C^{\mathrm{comp}} = C^{\mathrm{total}}_{\text{B1}} - C^{\mathrm{total}}_{\text{A1}}
\label{eq:competition_impact}
\end{equation}

\subsection{Interaction Effect}

The joint effect from combining REC and competition is

\begin{equation}
\Delta C^{\mathrm{int}} =
  \Bigl(C^{\mathrm{total}}_{\text{B2}} - C^{\mathrm{total}}_{\text{A1}}\Bigr)
  - \Delta C^{\mathrm{REC}} - \Delta C^{\mathrm{comp}}
\label{eq:interaction}
\end{equation}

\subsection{Battery Value}

The incremental value of battery optimisation relative to the REC-only baseline is

\begin{equation}
\Delta C^{\mathrm{batt}} = C^{\mathrm{total}}_{\text{C3}} - C^{\mathrm{total}}_{\text{A2}}
\label{eq:battery_value}
\end{equation}

A negative value indicates that the centrally optimised battery reduces total
supplier cost.

\subsection{Performance Indicators}

\paragraph{Forecast accuracy (MAPE).}

\begin{equation}
\mathrm{MAPE}_{n} = \frac{100\%}{T}
  \sum_{t \in \mathcal{T}}
  \left|\frac{\ell_{n}^{t} - \hat{\ell}_{n,\mathrm{ID}}^{t}}{\ell_{n}^{t}}\right|
\label{eq:mape}
\end{equation}

\paragraph{Relative imbalance exposure.}

\begin{equation}
R_{g}^{\mathrm{imb}} =
  \frac{\displaystyle\sum_{t \in \mathcal{T}} |\varepsilon_{g}^{t}|}
       {\displaystyle\sum_{t \in \mathcal{T}} |A_{g}^{t}|}
  \times 100\%
\label{eq:rel_imbalance}
\end{equation}

\paragraph{Average wholesale cost per kWh traded.}

\begin{equation}
\bar{c}_{s} =
  \frac{C_{s}^{\mathrm{total}}}
       {\displaystyle\sum_{t \in \mathcal{T}}\;\sum_{g \in \mathcal{G}_{s}} |A_{g}^{t}|}
\label{eq:avg_cost}
\end{equation}

%%=============================================================================
\section{Implementation Notes}
%%=============================================================================

\subsection{Temporal Resolution and Profile Scaling}

All physical quantities are read from 15-minute CSV files.  The mapping
between a normalised SimBench profile $\phi_{n}^{t}$ and the actual energy
volume is

\begin{equation}
\ell_{n}^{t} = \phi_{n}^{t} \cdot S_{n} \cdot \Delta t, \qquad
S_{n} = \frac{E_{n}^{\mathrm{target}}}{\sum_{t}\phi_{n}^{t}\,\Delta t}
\label{eq:profile_scaling}
\end{equation}

where $E_{n}^{\mathrm{target}}$ is the participant's annual energy target from
Table~1 in the profile mapping.

\subsection{Pipeline Activation per Scenario}

\begin{table}[H]
\centering
\caption{Market pipeline steps active in each scenario}
\label{tab:step_scope}
\begin{tabular}{@{}lccccc@{}}
\toprule
Step & A1 & A2 & B1 & B2 & C3\\
\midrule
i\;\; Day-Ahead Market        & $\checkmark$ & $\checkmark$ & $\checkmark$ & $\checkmark$ & $\checkmark$\\
ii\;\ Intra-Day Market        & $\checkmark$ & $\checkmark$ & $\checkmark$ & $\checkmark$ & $\checkmark$\\
ii-b Battery Optimisation     & --           & --           & --           & --           & $\checkmark$\\
iii\ REC Settlement           & --           & $\checkmark$ & --           & $\checkmark$ & $\checkmark$\\
iv\;\ Balancing Market        & $\checkmark$ & $\checkmark$ & $\checkmark$ & $\checkmark$ & $\checkmark$\\
v\;\;\ Supplier Billing       & $\checkmark$ & $\checkmark$ & $\checkmark$ & $\checkmark$ & $\checkmark$\\
\bottomrule
\end{tabular}
\end{table}

\subsection{Computational Complexity}

For $|\mathcal{S}|$ suppliers, $|\mathcal{G}|$ balancing groups, and
$|\mathcal{N}|$ participants, the simulation complexity of Steps~i--iii and~v
is $\mathcal{O}(T \cdot |\mathcal{N}| \cdot |\mathcal{G}|)$.  Step~ii-b
(MILP) adds $\mathcal{O}(T \cdot B)$ where $B = T = 35{,}136$ is the number
of binary variables.

\bibliography{references}
\bibliographystyle{plain}

\end{document}
